\chapter{\MakeUppercase{Introdução}}\label{chap:introducao}

O estudo das ciências naturais é frequentemente limitado pela capacidade humana de perceber o mundo pelos sentidos. Essas barreiras motivam a busca por alternativas que tornem os fenômenos mais acessíveis, facilitando a compreensão dos mecanismos que regem a natureza.

Um exemplo claro dessa limitação está na representação de moléculas e compostos químicos. Como ressaltado por Silva e Correia \cite{silva_ensino_2024}, a capacidade espacial de visualizar esses conceitos impacta diretamente sua compreensão, uma vez que tais estruturas são invisíveis a olho nu — o que dificulta o entendimento de noções como átomos, ligações químicas e interações moleculares. Além disso, a riqueza de detalhes de uma representação depende diretamente da habilidade de quem a produz. Representar moléculas longas, como polímeros, é uma tarefa tediosa; o desafio se intensifica ao considerar a natureza tridimensional dessas estruturas, que precisam ser retratadas em um plano bidimensional, tornando a representação inevitavelmente incompleta.

Apesar desses desafios, Rodrigo e Caridade \cite{rodrigues_historia_2024} observam que o advento da computação revolucionou as representações químicas, substituindo os símbolos tipográficos por visualizações ricas em cores e efeitos gráficos. A interseção entre química e computação se expande progressivamente, transformando o computador em um laboratório poderoso o suficiente para representar proteínas complexas e simular cenários com inúmeras variáveis. Ainda assim, a visualização permanece um fator central nos experimentos: como os dispositivos de exibição mais comuns são bidimensionais, mesmo que se simule um ambiente tridimensional, há um filtro imposto pela tela entre o usuário e a experiência. Esse limite é especialmente relevante para áreas como a cristalografia, que gera modelos 3D com detalhes atômicos de macromoléculas biológicas — essenciais, segundo Bazzano \cite{bazzano_modelagens_2022}, para a identificação de sítios ligantes de proteínas e para o desenvolvimento de fármacos.

Com o surgimento dos equipamentos de realidade estendida, abre-se a oportunidade de interagir com ambientes de realidade virtual. Na definição de Kirner \cite{kirner_claudio_evolucao_2011}, a realidade virtual é capaz de orquestrar um espaço tridimensional que, integrado aos sentidos humanos, cria possibilidades de interação antes inexistentes — superando a representação bidimensional convencional e ampliando os recursos de controle e exploração. Dessa convergência, espera-se que a abordagem supere os modelos tradicionais de ensino, acelerando tanto o aprendizado quanto o avanço científico.

\section{Descrição do Problema}\label{sec:descricao_do_problema}

A identificação visual de sítios ligantes em proteínas é uma etapa imprescindível no desenvolvimento de fármacos. Entretanto, os softwares de visualização molecular convencionais apresentam uma limitação estrutural: por operarem, majoritariamente, em dispositivos de apresentação planos, não oferecem ao usuário a sensação de profundidade necessária para compreender a geometria tridimensional do sítio de ligação, prejudicando a formação de uma noção espacial precisa da molécula.

Embora ferramentas com suporte à realidade virtual representem uma alternativa promissora a esse problema, as opções disponíveis atualmente são escassas e, em sua maioria, dependentes de hardware dedicado de alto desempenho gráfico — o que restringe o acesso a usuários com computadores de menor capacidade computacional. Ademais, essas soluções geralmente exigem instalação local, limitando ainda mais sua adoção em contextos acadêmicos e científicos com infraestrutura reduzida.

Diante desse cenário, identifica-se a necessidade de uma solução de visualização molecular em realidade virtual que seja leve, acessível e executável diretamente pelo navegador, sem dependência de hardware especializado, democratizando o acesso a essa forma de visualização.